\SGAchapter{XXV}{存在性定理}{Le théorème d’existence}
\begin{center}M. Demazure\end{center}
\SourcePage{267}
为求完整，本讲给出\Term{split-group}{分裂群}的存在性定理的一个证明。和 Chevalley 最初的证明（《关于某些\Term{semisimple-group-scheme}{半单群概形}》，Bourbaki 讨论班，1961 年 5 月，第 219 号）一样，这个证明以所有可能类型的复\Term{semisimple-group}{半单代数群}的存在性为基础。Cartier 曾给出一个更为令人满意的证明思路（未发表），直接证明与给定 \Term{cartan-matrix}{Cartan 矩阵}对应的\Term{simply-connected}{单连通}\Term{split-group}{分裂}\Term{semisimple-group}{半单} $\mathbf Z$-群存在。\footnote{编者注：B. Kostant [Ko66] 也曾勾勒出这样的思路；G. Lusztig 最近用量\Term{subgroup}{子群}给出了完整证明 [Lu09]。}不过应当指出，困难并不在于显式构造一个\Term{group-scheme}{群概形}，而在于验证所构造的群确实满足要求，也就是说，主要是验证它的\Term{fiber}{纤维}确实\Term{smooth}{光滑}且\Term{reductive-group}{约化}。

\SGAsection{1}{定理的陈述}
\begin{SGAresult}{定理}{1.1}
设 $S$ 为\Term{nonempty}{非空}\Term{scheme}{概形}。\Term{functor}{函子}
\[G\longmapsto\mathcal R(G)\]
给出\Term{pinning}{带钉扎}的\Term{reductive-group}{约化} $S$-\Term{group-scheme}{群概形}的\Term{category}{范畴}与\Term{pinning}{带钉扎}的\Term{reduced}{既约}\Term{root-datum}{根数据}的\Term{category}{范畴}之间的\Term{equivalence-category}{等价}。
\end{SGAresult}
根据唯一性定理（第 XXIII 讲 4.1），以上陈述\Term{equivalence}{等价}于：
\begin{SGAresult}{推论}{1.2（\Term{split-group}{分裂群}的存在性）}
对任意\Term{reduced}{既约}\Term{root-datum}{根数据} $\mathcal R$，存在\Term{reductive-group}{约化}\Term{split-group}{分裂} $\mathbf Z$-群 $G$，使得 $\mathcal R(G)\simeq\mathcal R$。\footnote{编者注：还应指出，F. Bruhat 与 J. Tits 的著作 [BT84] 包含 Chevalley 构造的一个变体（同处 2.2.3--2.2.5 及 \S 3.2），特别给出一个\Term{smooth}{光滑}\Term{affine}{仿射}、\Term{fiber}{纤维}\Term{connected}{连通}的 $\mathbf Z$-群 $G$，它具有\Term{split-torus}{分裂}\Term{maximal-torus}{极大环面}，且其\Term{generic-fiber}{一般纤维}是类型为 $\mathcal R$ 的\Term{reductive-group}{约化} $\mathbf Q$-群；$G$ 的\Term{geometric-fiber}{几何纤维}为\Term{reductive-group}{约化}的事实，可由同处 4.6.12 与 4.6.15 对特殊\Term{fiber}{纤维}的\Term{unipotent-radical}{幺幂根基}的研究推出（那些结果对由赋值\Term{root-datum}{根数据}所确定的更一般的 $G$ 也适用），但更简单的方法是由 $\mathfrak g=\operatorname{Lie}(G)$ 的描述、第 XIX 讲 1.12 (iii)，以及第 XX 讲 3.1 (iv) 中元素 $w_\alpha(X)$ 的存在性推出（参见 [BT84], 3.2.1）。}
\end{SGAresult}
\SourcePage{268}
特别地：
\begin{SGAresult}{推论}{1.3}
设 $k$ 为域。对于每个\Term{separable}{可分}裂的\Term{reductive-group}{约化} $k$-群 $G_k$，存在\Term{reductive-group}{约化} $\mathbf Z$-群 $G$，使得 $G\otimes_{\mathbf Z}k\simeq G_k$。
\end{SGAresult}
反过来，先注意，由第 XXII 讲 4.3.1 及第 XXI 讲 6.5.10，为证明 1.2，只需考虑\Term{root-datum}{根数据} $\mathcal R$ 为\Term{simply-connected}{单连通}的情形（如果愿意，根据第 XXI 讲 7.1.6，甚至可以假定它\Term{irreducible}{不可约}）。

另一方面，在 1.3 的条件下，\Term{group-scheme}{群概形} $G$ 的类型为常值（因为 $\operatorname{Spec}(\mathbf Z)$ \Term{connected}{连通}），因而其类型为 $\mathcal R(G_k)$；根据第 XXIII 讲 5.9，由此可知，对给定的群 $G_k$，1.3 成立即蕴涵存在类型为 $\mathcal R(G_k)$ 的\Term{split-group}{分裂} $\mathbf Z$-群。\footnote{编者注：事实上，无需使用第 XXIII 讲 5.9，因为本讲对每个\Term{semisimple-group}{半单} $\mathbf Q$-群 $G_{\mathbf Q}$，构造了与 $G_{\mathbf Q}$ 同类型并带有\Term{split-torus}{分裂}\Term{maximal-torus}{极大环面} $T$ 的\Term{semisimple-group}{半单} $\mathbf Z$-群 $G$；于是由第 XXII 讲 2.2，$G$ 为\Term{split-group}{分裂}的。}

因此，为证明 1.2，进而证明 1.1，只需在 $k$ 的\Term{characteristic}{特征}为零（例如 $k=\mathbf C$）且 $G_k$ \Term{simply-connected}{单连通}（特别是\Term{semisimple-group}{半单}）时证明 1.3，并证明：
\begin{SGAresult}{命题}{1.4}
对于每个\Term{reduced}{既约}\Term{simply-connected}{单连通}\Term{root-datum}{根数据} $\mathcal R$，存在类型为 $\mathcal R$ 的\Term{semisimple-group}{半单} $\mathbf C$-\Term{algebraic-group}{代数群}。
\end{SGAresult}
可以用如下方法证明 1.4。首先，已知存在类型为 $\mathcal R$ 的复\Term{semisimple-lie-algebra}{半单 Lie 代数} $\mathfrak g$，例如参见 N. Jacobson, \emph{Lie Algebras}, 第 VII 章，定理 5。\footnote{编者注：设 $\Delta$ 为 $R$ 的一组基，$\widetilde{\mathfrak g}$ 为由\Term{generator}{生成元} $(e_\alpha,f_\alpha,h_\alpha)_{\alpha\in\Delta}$ 及关系
\[
[h_\alpha,h_\beta]=0,\qquad
[e_\alpha,f_\beta]=\begin{cases}h_\alpha&\beta=\alpha,\\0&\beta\ne\alpha,\end{cases}
\]
\[
[h_\alpha,e_\beta]=(\alpha^*,\beta)e_\beta,\qquad
[h_\alpha,f_\beta]=-(\alpha^*,\beta)f_\beta
\]
定义的 $\mathbf C$-\Term{lie-algebra}{Lie 代数}。在所引文献中，$\mathfrak g$ 定义为 $\widetilde{\mathfrak g}$ 对其所有\Term{finite}{有限}维\Term{irreducible}{不可约}表示的核之交的商；[Se66], \S VI.5，定理 9（也参见 [BLie], VIII, \S 4.3，定理 1）证明了 $\mathfrak g$ 是 $\widetilde{\mathfrak g}$ 对关系
\[
\operatorname{ad}(e_\alpha)^{1-(\alpha^*,\beta)}(e_\beta)=0,
\qquad
\operatorname{ad}(f_\alpha)^{1-(\alpha^*,\beta)}(f_\beta)=0
\]
所得的商。结构常数的显式描述（尤其是\Term{signature}{符号}的选取）见第 XXIII 讲 6.5 和 6.7，以及 [Ti66], \S 4，定理 1；对于类型 $A,D,E$，另见 [Sp98], 10.2.5。}于是 $G=\operatorname{Aut}(\mathfrak g)^0$ 是类型为 $\operatorname{ad}(\mathcal R)$ 的\Term{semisimple-group}{半单} $\mathbf C$-\Term{algebraic-group}{代数群}。\footnote{编者注：可以假定 $\mathcal R$ \Term{irreducible}{不可约}，从而 $\mathfrak g$ 为单的。由一般论证，$\operatorname{Lie}(G)$ 是 $\mathfrak g$ 的\Term{derivation}{导子}所成的 $\mathbf C$-\Term{lie-algebra}{Lie 代数}（参见 [DG70], \S II.4，命题 2.3），而这些\Term{derivation}{导子}全是内\Term{derivation}{导子}（参见 [BLie], I, \S 6.1，命题 1 的推论 3），故 $\operatorname{Lie}(G)=\mathfrak g$；于是 $G$ 没有\Term{dimension}{维数} $>0$ 的\Term{invariant-subgroup}{不变子群}，因而 $G$ 为\Term{semisimple-group}{半单}的。它的\Term{root-system}{根系}于是与 $\mathfrak g$ 的\Term{root-system}{根系}相同；此外，$G$ 的\Term{center}{中心}在 $\mathfrak g$ 上的\Term{action}{作用}既平凡又忠实，故该\Term{center}{中心}平凡，因而 $G$ 是\Term{adjoint-group}{伴随型}的。}根据 Bible, \S 23.1，命题 1，即可推出类型为 $\mathcal R$ 的\Term{semisimple-group}{半单} $\mathbf C$-群存在。

本讲余下的内容用来在 $k$ 的\Term{characteristic}{特征}为零且 $G_k$ \Term{semisimple-group}{半单}时证明 1.3。证明分两步：构造一个“$\mathbf Z$-\Term{group-chunk}{群概形片段}”（第 2 节），再研究应用“Weil 定理”所得的群（第 3 节）。

为避免混淆，在第 2 节中，若 $X$ 是 $\mathbf Z$-\Term{scheme}{概形}，我们不使用缩写 $X_k$ 表示 $k$-\Term{scheme}{概形} $X\otimes_{\mathbf Z}k$。
\SourcePage{269}
\SGAsection{2}{存在性定理：群片段的构造}
\textbf{2.1.} 一次性选定 $G_k$ 的一个\Term{splitting}{分裂结构}，记为
\[(G_k,T_k,M,R)\]
（参见第 XXII 讲 1.13），选定 $R$ 的一个\Term{simple-root}{单根}系 $\Delta$（它确定\Term{positive-root}{正根}系 $R_+$），并选定 $G_k$ 的一个 \Term{chevalley-system}{Chevalley 系} $(X_{\alpha,k})_{\alpha\in R}$（第 XXIII 讲 6.1 和 6.2），使之满足如下附加条件（参见第 XX 讲 2.6）：对所有 $\alpha\in R$，有
\[X_{\alpha,k}X_{-\alpha,k}=1.\]
最后，在 $M$ 中由 $R$ 生成的\Term{subgroup}{子群}上选定一个与\Term{group-structure}{群结构}相容的全\Term{order-relation}{序关系}，使得大于零的根恰好是 $R_+$ 的元素。于是将各根记为
\[-\alpha_n<-\alpha_{n-1}<\cdots<-\alpha_1<\alpha_1<\alpha_2<\cdots<\alpha_n.\]
对 $\alpha\in R$，用 $U_{\alpha,k}$ 表示对应于根 $\alpha$ 的\Term{vector-group}{向量群}，用
\[p_{\alpha,k}:\mathbf G_{\mathrm a,k}\xrightarrow{\sim}U_{\alpha,k}\]
表示由 $X_{\alpha,k}$ 定义的\Term{vector-group}{向量群}\Term{isomorphism}{同构}。

\textbf{2.2.} $G_k$ 的\Term{splitting}{分裂结构}特别包含一个 $k$-群\Term{isomorphism}{同构}
\[T_k\simeq\operatorname D_k(M).\]
令 $T=\operatorname D(M)$；它是一个 $\mathbf Z$-\Term{torus}{环面}，并且可将以上\Term{isomorphism}{同构}视为一个\Term{isomorphism}{同构}
\[T_k\simeq T\otimes_{\mathbf Z}k.\]
有
\[\operatorname{Hom}_{\mathbf Z\text{-群}}(T,\mathbf G_{\mathrm m})=M,\]
我们将 $R\subset M$ 中的元素视为 $T$ 的\Term{character}{特征标}。同样，将 $R^*$ 中的元素视为 $\mathbf Z$-群\Term{morphism}{态射} $\mathbf G_{\mathrm m}\to T$。

\textbf{2.3.} 对每个 $\alpha\in R_+$，取 $\mathbf G_{\mathrm a}$ 的一个副本 $\mathbf G_{\mathrm a}(\alpha)$；考虑 $\mathbf Z$-\Term{scheme}{概形}
\[U=\mathbf G_{\mathrm a}(\alpha_1)\times\cdots\times\mathbf G_{\mathrm a}(\alpha_n).\]
设 $U_k$ 表示 $G_k$ 中由 $R_+$ 定义的 \Term{borel-subgroup}{Borel 子群} $B_k$ 的\Term{unipotent}{幺幂}部分，用
\[a:U\otimes_{\mathbf Z}k\xrightarrow{\sim}U_k\]
表示由
\[a(x_1,\ldots,x_n)=p_{\alpha_1,k}(x_1)\cdots p_{\alpha_n,k}(x_n)\]
定义的 $k$-\Term{scheme}{概形}\Term{isomorphism}{同构}。

\SourcePage{270}
\textbf{2.4.} $U_k$ 的群律写成如下\Term{form}{形式}的关系：
\[a(x_1,\ldots,x_n)\cdot a(y_1,\ldots,y_n)=a(z_1,\ldots,z_n),\]
其中每个 $z_h$（$h=1,\ldots,n$）均表达为多项式
\[z_h=x_h+y_h+Q_h(x_1,\ldots,x_{h-1},y_1,\ldots,y_{h-1}),\]
$Q_h$ 的系数为整数（第 XXII 讲 5.5.8 及第 XXIII 讲 6.4）。此外，
\[Q_h(x_1,\ldots,x_{h-1},0,\ldots,0)=0.\]
赋予 $U$ 由以上公式定义的合成律（这些公式确实“定义在 $\mathbf Z$ 上”）。因为它在 $U_k$ 上诱导其群律，故该合成律满足结合律，而且 $(0)$ 是单位元（事实上，以上两条断言均可用多项式 $Q_h$ 之间的关系表达，而 $\mathbf Z\to k$ 是\Term{injectivity}{单射}）。下面说明这是一条群律：若 $(x_i)$ 是 $U$ 的一个\Term{section}{截面}（在任意 $S$ 上），则可由递推公式
\[y_i=-x_i-Q_i(x_i,\ldots,x_{i-1},y_1,\ldots,y_{i-1})\]
计算 $(x_i)$ 的逆元 $(y_i)$；这些公式也“定义在 $\mathbf Z$ 上”。

综上，我们在 $U$ 上构造了一条群律，使得以上\Term{isomorphism}{同构} $a$ 成为群\Term{isomorphism}{同构}。

对每个 $\alpha\in R_+$，考虑\Term{morphism}{态射}
\[p_\alpha:\mathbf G_{\mathrm a}\longrightarrow U,\]
其定义为 $p_\alpha(x)=(x_i)$，其中
\[x_i=\begin{cases}x&\text{若 }\alpha_i=\alpha,\\0&\text{若 }\alpha_i\ne\alpha.\end{cases}\]
它是一个\Term{closed-immersion}{闭浸入}，也是\Term{group-homomorphism}{群同态}；用 $U_\alpha$ 表示其像。有 $(x_i)=p_{\alpha_1}(x_1)\cdots p_{\alpha_n}(x_n)$，这说明 $U$ 可与如下乘积等同：
\[U=U_{\alpha_1}\cdot U_{\alpha_2}\cdots U_{\alpha_n}.\]

\textbf{2.5.} 使 $T=\operatorname D(M)$ 通过\Term{character}{特征标} $\alpha$ \Term{action}{作用}于每个 $U_\alpha$；立即可验证，这定义了 $T$ 在群 $U$ 上的\Term{action}{作用}，因而\Term{constructible}{可构造}\Term{semidirect-product}{半直积} $B=T\cdot U$。有一个\Term{canonical}{典范}的 $k$-群\Term{isomorphism}{同构}
\[B\otimes_{\mathbf Z}k\xrightarrow{\sim}B_k.\]
现在，若在 $R_+$ 上取任意一个次序，则由 $U$ 中的乘法定义的\Term{morphism}{态射}
\[\prod_{\alpha\in R_+}U_\alpha\longrightarrow U\]
仍是\Term{isomorphism}{同构}。事实上，因为两边都是\Term{flat}{平坦}且\Term{finite-presentation}{有限呈示}的 $\mathbf Z$-\Term{scheme}{概形}，只需在\Term{geometric-fiber}{几何纤维}上验证这一断言；于是归结为 Lazard 理论（Bible, \S 13.1，命题 1）：将 $U$ 看作以 $T$ 为\Term{operator}{算子}群的群，并利用各 $U_\alpha$ 作为带\Term{operator}{算子}的群两两不\Term{isomorphism}{同构}这一事实（因为 $T$ 的\Term{character}{特征标} $\alpha\in R_+$ 在每个\Term{fiber}{纤维}上两两不同）。

\SourcePage{271}
\textbf{2.6.} 将 $R_+$ 换为 $R_-=-R_+$，可同样构造群 $U^-,B^-,U_\alpha$（$\alpha\in R_-$）及\Term{isomorphism}{同构}
\[p_\alpha:\mathbf G_{\mathrm a}\xrightarrow{\sim}U_\alpha,\qquad\alpha\in R_-.\]
最后引入乘积\Term{scheme}{概形}
\[\Omega=U^-\times T\times U;\]
有一个\Term{canonical}{典范}的 $k$-\Term{scheme}{概形}\Term{isomorphism}{同构}
\[\Omega\otimes_{\mathbf Z}k\xrightarrow{\sim}U_k^-\times_kT_k\times_kU_k\simeq\Omega_k,\]
其中 $\Omega_k$ 是 $G_k$ 的“\Term{big-cell}{大胞腔}”（第 XXII 讲 4.1.2）。

从现在起，通过以上\Term{isomorphism}{同构}将 $\Omega\otimes_{\mathbf Z}k$ 与 $\Omega_k$ 等同；借助\Term{unit-section}{单位截面}，将 $U^-,T,U$ 视为 $\Omega$ 的子\Term{scheme}{概形}。用 $\underline e=((0),e,(0))$ 表示 $\Omega$ 的“\Term{unit-section}{单位截面}”。

现在的目标是在 $\Omega$ 上给出一条\Term{group-chunk}{群片段}律。
\begin{SGAresult}{引理}{2.7}
设 $\alpha\in\Delta$，$w_{\alpha,k}$ 为由 $X_{\alpha,k}$ 定义的 $\operatorname{Norm}_{G_k}(T_k)(k)$ 中的元素（回忆其定义为
\[w_{\alpha,k}=p_{-\alpha,k}(-1)p_{\alpha,k}(1)p_{-\alpha,k}(-1)\]）。
存在 $\Omega$ 的一个包含\Term{section}{截面} $\underline e$ 的\Term{open-subset}{开集} $V_\alpha$，以及\Term{morphism}{态射}
\[h_\alpha:V_\alpha\longrightarrow\Omega,\]
满足以下条件：

(i) $h_\alpha(\underline e)=\underline e$。

(ii) $(h_\alpha)\otimes_{\mathbf Z}k$ 与 $\operatorname{int}(w_{\alpha,k})$ 在 $V_\alpha\otimes_{\mathbf Z}k\subset G_k$ 上的\Term{restriction}{限制}相同。

(iii) $T\subset V_\alpha$，且 $h_\alpha$ 将 $T$ 映入 $T$。对每个 $\beta\in R$，有 $U_\beta\subset V_\alpha$，且 $h_\alpha$ 将 $U_\beta$ 映入 $U_{s_\alpha(\beta)}$。
\end{SGAresult}
\begin{proof}
根据 \Term{chevalley-system}{Chevalley 系}的定义（第 XXIII 讲 6.1），对每个 $\beta\in R$，存在整数 $e_\beta=\pm1$，使对所有 $x\in\mathbf G_{\mathrm a}(S)$，$S\to\operatorname{Spec}(k)$，有
\[\operatorname{int}(w_{\alpha,k})p_{\beta,k}(x)=p_{s_\alpha(\beta),k}(e_\beta x).\]
设 $S$ 为任意\Term{scheme}{概形}，将 $\Omega(S)$ 的任意元素写成
\[
u=\left(
 \left(\prod_{\substack{\beta\in R_-\\\beta\ne-\alpha}}p_\beta(x_\beta)\right)p_{-\alpha}(x_{-\alpha}),\quad
 t,\quad
 p_\alpha(x_\alpha)\left(\prod_{\substack{\beta\in R_+\\\beta\ne\alpha}}p_\beta(x_\beta)\right)
\right),
\]
其中在 $R_-\setminus\{-\alpha\}$ 与 $R_+\setminus\{\alpha\}$ 上选定了（任意的）次序（参见 2.5）。定义\Term{morphism}{态射} $d:\Omega\to\operatorname{Spec}(\mathbf Z)$ 为
\[d(u)=\alpha(t)+x_\alpha x_{-\alpha}.\]
\SourcePage{272}
设 $V_\alpha$ 为\Term{open-subset}{开集} $\Omega_d$（即 $\Omega$ 中由“$d(u)$ \Term{invertible}{可逆}”定义的\Term{open-subset}{开集}）；它包含 $\underline e,T$ 以及每个 $U_\beta$，$\beta\in R$。令
\[h_\alpha:V_\alpha\longrightarrow\Omega\]
为由 $h_\alpha(u)=(a(u),b(u),c(u))$ 定义的\Term{morphism}{态射}，其中
\begin{align*}
a(u)&=\left(\prod_{\substack{\beta\in R_-\\\beta\ne-\alpha}}p_{s_\alpha(\beta)}(e_\beta x_\beta)\right)\cdot p_{-\alpha}(-x_\alpha d(u)^{-1}),\\
b(u)&=t\cdot\alpha^*(d(u)),\\
c(u)&=p_\alpha(-x_{-\alpha}d(u)^{-1})\cdot\left(\prod_{\substack{\beta\in R_+\\\beta\ne\alpha}}p_{s_\alpha(\beta)}(e_\beta x_\beta)\right).
\end{align*}
由于 $s_\alpha$ 置换异于 $\alpha$ 的\Term{positive-root}{正根}（相应地，异于 $-\alpha$ 的\Term{negative-root}{负根}），$c(u)$（相应地，$a(u)$）是 $U$（相应地，$U^-$）的\Term{section}{截面}，所以上述\Term{morphism}{态射}确有定义。它显然满足 (i) 与 (iii)。至于 (ii)，则立即由 $e_\beta$（$\beta\in R$）的定义及第 XX 讲 3.12 得出。
\end{proof}

\begin{SGAresult}{引理}{2.8}
存在 $\Omega$ 的\Term{open-subset}{开集} $V,V'$ 及\Term{morphism}{态射}
\[h:V\longrightarrow\Omega,\qquad h':V'\longrightarrow\Omega,\]
满足以下条件：

(i) $V,V'$ 包含 $\underline e$，且 $h(\underline e)=h'(\underline e)=\underline e$。

(ii) $h'\circ h:h^{-1}(V')\to\Omega$ 诱导的\Term{morphism}{态射}是恒等\Term{morphism}{态射}的\Term{restriction}{限制}。

(iii) $k$-\Term{morphism}{态射} $h\otimes_{\mathbf Z}k$ 与 $h'\otimes_{\mathbf Z}k$ 是群 $G_k$ 的\Term{automorphism}{自同构}分别在 $V\otimes_{\mathbf Z}k$ 与 $V'\otimes_{\mathbf Z}k$ 上的\Term{restriction}{限制}。

(iv) $V,V'$ 包含 $U,T,U^-$；$h,h'$ 将 $U$ 映入 $U^-$，将 $U^-$ 映入 $U$，并将 $T$ 映入 $T$。
\end{SGAresult}
\begin{proof}
设 $\overline w_0$ 为 $G_k$ 的 \Term{weyl-group}{Weyl 群}中将 $R_+$ 变为 $R_-$ 的元素。写成
\[\overline w_0=s_{\alpha_n}\cdots s_{\alpha_1},\qquad\alpha_i\in\Delta\]
（此处的编号与 2.1 无关）。令
\[w_0=w_{\alpha_n,k}\cdots w_{\alpha_1,k}\in\operatorname{Norm}_{G_k}(T_k)(k).\]
对 $i\leq n$，递归定义 $\Omega$ 的\Term{open-subset}{开集} $V_i$ 及\Term{morphism}{态射} $h_i:V_i\to\Omega$：令 $V_0=\Omega$，$h_0=\operatorname{id}$，并对 $i=0,\ldots,n-1$ 令
\[V_{i+1}=h_i^{-1}(V_{\alpha_{i+1}}),\qquad h_{i+1}=h_{\alpha_{i+1}}\circ h_i,\]
其中 $V_{\alpha_j}$ 与 $h_{\alpha_j}$ 为 2.7 中的记号。
\SourcePage{273}
取 $V=V_n$，$h=h_n$。关于 $V,h$ 的条件 (i)、(iii)、(iv) 确实成立；对于 (i) 与 (iii)，这立即由 2.8 得出；对于 (iv)，则因为 $h\otimes_{\mathbf Z}k$ 是 $\operatorname{int}(w_0)$ 在 $V\otimes_{\mathbf Z}k$ 上的\Term{restriction}{限制}。

由于 $(\overline w_0)^2=1$，还有
\[\overline w_0=s_{\alpha_1}\cdots s_{\alpha_n}=(s_{\alpha_1}^{3})\cdots(s_{\alpha_n}^{3}).\]
令
\[w'_0=(w_{\alpha_1,k})^3\cdots(w_{\alpha_n,k})^3,\]
作与上述相同的构造，即得一个\Term{open-subset}{开集} $V'$ 及\Term{morphism}{态射} $h'$，它们也满足 (i)、(iii)、(iv)。此外，$h'\otimes_{\mathbf Z}k$ 是 $\operatorname{int}(w'_0)$ 在 $V'\otimes_{\mathbf Z}k$ 上的\Term{restriction}{限制}。但是，对每个\Term{simple-root}{单根} $\alpha\in\Delta$，有 $(w_{\alpha,k})^4=e$（参见第 XX 讲 3.1），故 $w'_0w_0=e$。这说明 $h'\circ h$ 在 $\Omega\otimes_{\mathbf Z}k$ 的某个\Term{nonempty}{非空}\Term{open-subset}{开集}中诱导恒等\Term{morphism}{态射}。而由于 $\Omega$ 在 $\mathbf Z$ 上\Term{smooth}{光滑}且\Term{finite-presentation}{有限呈示}，$\Omega\otimes_{\mathbf Z}k$ 在 $\Omega$ 中\Term{schematically-dense}{概形稠密}，从而证明了 (ii)。
\end{proof}

\begin{SGAresult}{命题}{2.9}
存在 $\Omega\times\Omega$ 的\Term{open-subset}{开集} $V_1$、$\Omega$ 的\Term{open-subset}{开集} $V_2$ 以及\Term{morphism}{态射}
\[\pi:V_1\longrightarrow\Omega,\qquad\sigma:V_2\longrightarrow\Omega,\]
具有以下性质：

(i) 若 $x\in\Omega(S)$，则 $(\underline e,x)$ 与 $(x,\underline e)$ 是 $V_1$ 的\Term{section}{截面}，且
\[\pi(\underline e,x)=\pi(x,\underline e)=x.\]
(ii) $V_2$ 包含 $\underline e$，且 $\sigma(\underline e)=\underline e$。

(iii) $\pi_k$ 与 $\sigma_k$ 是由乘法（相应地，取逆）定义的\Term{morphism}{态射} $G_k\times_kG_k\to G_k$ 及 $G_k\to G_k$ 的\Term{restriction}{限制}。
\end{SGAresult}
\begin{proof}
设 $(v,t,u)$ 与 $(v',t',u')$ 是 $\Omega$ 的两个\Term{section}{截面}。由 2.8 (iv)，$h(u)$ 是 $U^-$ 的\Term{section}{截面}，$h(v')$ 是 $U$ 的\Term{section}{截面}，从而可将 $(h(u),e,h(v'))$ 视为 $\Omega$ 的\Term{section}{截面}。设 $V_1$ 为由条件
\[((v,t,u),(v',t',u'))\in V_1(S)\quad\Longleftrightarrow\quad(h(u),e,h(v'))\in V'(S)\]
定义的 $\Omega\times\Omega$ 的\Term{open-subset}{开集}（记号见 2.8）。若 $((v,t,u),(v',t',u'))$ 是 $V_1$ 的\Term{section}{截面}，则 $h'(h(u),e,h(v'))$ 有定义；它是 $\Omega$ 的\Term{section}{截面}，\Term{decomposable}{可分解}为
\[h'(h(u),e,h(v'))=(v'',t'',u'').\]
于是令
\[\pi((v,t,u),(v',t',u'))=(v\cdot tv''t^{-1},\;tt''t',\;t'^{-1}u''t'\cdot u').\]
(i) 的验证是直接的（由 2.8 (ii)）。为验证 (iii) 中关于 $\pi$ 的条件，注意到当 $u\in U(S)$、$v\in U^-(S)$、$S\to k$ 时，由 2.8 (iii)、(ii)，有
\[h'(h(u),e,h(v'))=uv'=v''t''u''.\]
$\sigma$ 的构造类似：若 $(v,t,u)$ 是 $\Omega$ 的\Term{section}{截面}，则 $h(u^{-1})$ 是 $U^-$ 的\Term{section}{截面}，
\SourcePage{274}
$h(v^{-1})$ 是 $U$ 的\Term{section}{截面}，$h(t^{-1})$ 是 $T$ 的\Term{section}{截面}，所以 $(h(u^{-1}),h(t^{-1}),h(v^{-1}))$ 是 $\Omega$ 的\Term{section}{截面}。于是可用
\[(v,t,u)\in V_2(S)\quad\Longleftrightarrow\quad(h(u^{-1}),h(t^{-1}),h(v^{-1}))\in V'(S)\]
定义 $\Omega$ 的\Term{open-subset}{开集} $V_2$，并用
\[\sigma(v,t,u)=h'(h(u^{-1}),h(t^{-1}),h(v^{-1}))\]
定义\Term{morphism}{态射} $\sigma:V_2\to\Omega$。关于 $\sigma$ 的条件可如上验证。
\end{proof}

\begin{SGAresult}{推论}{2.10}
$\pi$ “一般地满足结合律”，而 $\sigma$ 是“一般的逆”：若 $x,y,z\in\Omega(S)$ 且以下表达式有定义（在 $\Omega$ 的一个包含\Term{unit-section}{单位截面}的\Term{open-subset}{开集}上，这总会发生），则有
\[\pi(x,\pi(y,z))=\pi(\pi(x,y),z),\qquad\pi(x,\sigma(x))=\underline e=\pi(\sigma(x),x).\]
\end{SGAresult}
事实上，每个公式的两边都定义了\Term{smooth}{光滑}且\Term{finite-presentation}{有限呈示}的 $\mathbf Z$-\Term{scheme}{概形}之间的\Term{morphism}{态射}；由 2.10 (iii)，它们在\Term{generic-fiber}{一般纤维}上相同。

\begin{SGAresult}{推论}{2.11}
设 $\alpha\in R$。对任意 $S$ 及任意满足 $(p_\alpha(x),p_{-\alpha}(y))\in V_1(S)$ 且 $1+xy\in\mathbf G_{\mathrm m}(S)$ 的 $x,y\in\mathbf G_{\mathrm a}(S)$（这些条件定义了 $\mathbf G_{\mathrm a,S}^2$ 中包含\Term{section}{截面} $(0,0)$ 的\Term{open-subset}{开集}），有
\[\pi(p_\alpha(x),p_{-\alpha}(y))=
\left(p_{-\alpha}\left(\frac{y}{1+xy}\right),\;\alpha^*(1+xy),\;p_\alpha\left(\frac{x}{1+xy}\right)\right).\]
\end{SGAresult}
根据第 XX 讲 2.1，证明与前面的证明相同。

\SGAsection{3}{存在性定理：证明的完成}
为简化措辞，先作如下定义。
\begin{SGAresult}{定义}{3.1}
设 $S$ 为\Term{scheme}{概形}，$G$ 为 $S$-\Term{group-scheme}{群概形}。若存在 $S$-\Term{scheme}{概形}的\Term{open-immersion}{开浸入} $i:\Omega_S=\Omega\times S\to G$，满足以下条件，则称 $G$ \Term{admissible-group}{容许}：

(i) 图
\[
\begin{tikzcd}
(V_1)_S\arrow[r,hook]\arrow[d,"\pi"']&\Omega_S\times_S\Omega_S\arrow[r,"i\times_Si"]&G\times_SG\arrow[d,"\pi_G"]\\
\Omega_S\arrow[rr,"i_S"']&&G
\end{tikzcd}
\]
交换，其中 $\pi_G$ 表示 $G$ 中的乘法\Term{morphism}{态射}。

(ii) 存在\Term{finite}{有限}个\Term{section}{截面} $a_j\in\Omega(S)$，使各 $i(a_j)\cdot i(\Omega_S)$ \Term{covering}{覆盖} $G$。
\end{SGAresult}
由“Weil 定理”（第 XVIII 讲 3.13 (iii)、(iv)），有：
\begin{SGAresult}{引理}{3.2}
若对每个在 $\mathbf Z$ 上\Term{etale}{平展}且\Term{finite-type}{有限型}的\Term{scheme}{概形} $S$，每个\Term{admissible-group}{容许} $S$-群都\Term{affine}{仿射}，则存在一个\Term{admissible-group}{容许}且\Term{affine}{仿射}的 $\mathbf Z$-群。
\end{SGAresult}

\SourcePage{275}
现在有：
\begin{SGAresult}{引理}{3.3}
设 $S$ 为\Term{scheme}{概形}，$G$ 为\Term{admissible-group}{容许} $S$-群。则 $G$ 在 $S$ 上\Term{smooth}{光滑}且\Term{finite-presentation}{有限呈示}，其\Term{fiber}{纤维}\Term{affine}{仿射}、\Term{connected}{连通}且\Term{semisimple-group}{半单}。
\end{SGAresult}
\begin{proof}
因为 $\Omega_S$ 在 $S$ 上\Term{smooth}{光滑}且\Term{finite-presentation}{有限呈示}，并具有\Term{connected}{连通}\Term{fiber}{纤维}，由条件 (ii)，$G$ 也如此。因此，为验证 3.3，可假定 $S$ 是域 $K$ 的谱。将 $\Omega_K$ 与它在 $G$ 中的像等同。显然，$\Omega_K$ 是 $G$ 的\Term{subgroup}{子群} $T_K$ 及 $U_{\alpha,K}$（$\alpha\in R$）的乘积
\[\prod_{\alpha\in R_-}U_{\alpha,K}\cdot T_K\cdot\prod_{\alpha\in R_+}U_{\alpha,K}.\]
因此 $G$ 的\Term{lie-algebra}{Lie 代数}可与\Term{direct-sum}{直和}
\[\operatorname{Lie}(T_K)\oplus\bigoplus_{\alpha\in R}\operatorname{Lie}(U_{\alpha,K})\]
等同。因为由 $T_K$ 的一个\Term{section}{截面}定义的\Term{inner-automorphism}{内自同构}通过\Term{character}{特征标}
\[\alpha\in R\subset M\simeq\operatorname{Hom}_{K\text{-群}}(T_K,\mathbf G_{\mathrm m,K})\]
\Term{action}{作用}于 $U_{\alpha,K}$，进而\Term{action}{作用}于 $\operatorname{Lie}(U_{\alpha,K})$，所以 $\operatorname{Lie}(G)$ 的上述\Term{resolution}{分解}恰是 $T$ 的\Term{adjoint-action}{伴随作用}下的\Term{resolution}{分解}。因此 $G_K$ 相对于 $T_K$ 的根正是各 $\alpha\in R$。

应用第 XIX 讲 1.13。设 $T_\alpha$ 为 $\operatorname{Ker}(\alpha)\subset T_K$ 的\Term{maximal-torus}{极大环面}，并设 $Z_\alpha=\operatorname{Centr}_G(T_\alpha)$；只需证明每个 $Z_\alpha$ 为\Term{reductive-group}{约化}的。而 $Z_\alpha\cap\Omega_K$ 正是
\[\prod_{\substack{\beta\in R_-\\\beta|_{T_\alpha}=e}}U_{\beta,K}\cdot T_K\cdot\prod_{\substack{\beta\in R_+\\\beta|_{T_\alpha}=e}}U_{\beta,K};\]
但在 $T_\alpha$ 上为零的根是 $\alpha$ 的有理倍数，也就是 $\alpha$ 和 $-\alpha$；这说明
\[Z_\alpha\cap\Omega_K=U_{-\alpha,K}\cdot T_K\cdot U_{\alpha,K}.\]
根据第 XX 讲 3.4，为证明 $Z_\alpha$ 为\Term{reductive-group}{约化}的，只需证明 $U_{\alpha,K}$ 与 $U_{-\alpha,K}$ 不交换，而这立即由 2.11 得到。
\end{proof}
由 3.2 和 3.3，只要证明以下引理，整个证明就完成了。
\begin{SGAresult}{引理}{3.4}
若 $S$ 为\Term{dimension}{维数} $\leq1$ 的\Term{locally-noetherian}{局部 Noether}\Term{scheme}{概形}，$G$ 为\Term{smooth}{光滑}\Term{finite-type}{有限型} $S$-群，且其\Term{fiber}{纤维}\Term{affine}{仿射}、\Term{connected}{连通}、\Term{semisimple-group}{半单}，则 $G$ \Term{affine}{仿射}（从而\Term{semisimple-group}{半单}）。
\end{SGAresult}
\textbf{说明。} 在第 XVI 讲中已经看到，3.4 无需对 $S$ 作任何假设也成立，但其证明相当精细；因为此处只需特例 3.4，我们给出一个直接证明。
\begin{proof}
考虑 $G$ 的\Term{lie-algebra}{Lie 代数} $\mathfrak g$，它是\Term{locally-free}{局部自由}的 $\mathcal O_S$-模，并考虑 $G$ 的\Term{adjoint-representation}{伴随表示}
\[\operatorname{Ad}:G\longrightarrow\operatorname{GL}_{\mathcal O_S}(\mathfrak g).\]
\SourcePage{276}
为证明 $G$ 在 $S$ 上\Term{affine}{仿射}，只需证明\Term{morphism}{态射} $\operatorname{Ad}$ \Term{affine}{仿射}。因为 $G$ \Term{smooth}{光滑}且\Term{fiber}{纤维}\Term{connected}{连通}，它在 $S$ 上\Term{separated}{分离}（第 VI$_{\mathrm B}$ 讲 5.5），故\Term{morphism}{态射} $\operatorname{Ad}$ \Term{separated}{分离}。利用附录中证明的一个结果（见 4.1），只需证明\Term{morphism}{态射} $\operatorname{Ad}$ \Term{quasi-finite}{拟有限}。于是归结为 $S$ 是域的谱的情形；此时 $G$ \Term{affine}{仿射}，因而\Term{semisimple-group}{半单}，所以归结为第 XXII 讲 5.7.14。
\end{proof}

\SGAsection{4}{附录}
我们在证明过程中使用了以下命题。
\begin{SGAresult}{命题}{4.1}
设 $S$ 为\Term{dimension}{维数} $\leq1$ 的\Term{locally-noetherian}{局部 Noether}\Term{scheme}{概形}，$G,H$ 为两个\Term{finite-type}{有限型} $S$-\Term{group-scheme}{群概形}，$f:G\to H$ 为\Term{quasi-finite}{拟有限}且\Term{separated}{分离}的群\Term{morphism}{态射}。若 $G$ 在 $S$ 上\Term{flat}{平坦}，\footnote{编者注：我们补入了原先遗漏的\Term{flat}{平坦}性假设。}则 $f$ \Term{affine}{仿射}。
\end{SGAresult}
我们只在 $G$ 在 $S$ 上\Term{smooth}{光滑}时给出证明；在我们使用该命题的情形中，这个假设确实满足。

\textbf{4.2.} 根据 EGA II, 1.6.4，可以假定 $S$ \Term{reduced}{既约}。通过通常的取\Term{limit}{极限}技巧，\footnote{编者注：参见 EGA IV$_3$, 8.10.5 (viii)。}可以假定 $S$ 为\Term{local-scheme}{局部}的。若 $\dim(S)=0$，断言是直接的；\footnote{编者注：事实上，若 $S=\operatorname{Spec}(k)$（$k$ 为域），则 $f$ 是\Term{projection}{投影} $p:G\to G/\operatorname{Ker}(f)$ 与一个\Term{closed-immersion}{闭浸入} $i$ 的复合；因 $\operatorname{Ker}(f)$ 在 $k$ 上\Term{finite}{有限}，故 $p$ \Term{finite}{有限}（第 VI$_{\mathrm B}$ 讲 9.2）。}现假设 $\dim(S)=1$。由\Term{faithfully-flat-descent}{忠实平坦下降}，可以假定 $S$ \Term{complete}{完备}且\Term{residue-field}{剩余域}\Term{algebraically-closed}{代数闭}。将 $S$ 换为其\Term{normalization}{正规化} $\widetilde S$ 后，可以假定 $S$ \Term{normal}{正规}（EGA II, 6.7.1 及 EGA 0$_{\mathrm{IV}}$, 23.1.5）。\footnote{编者注：事实上，$S$ 的\Term{irreducible-component}{不可约分支} $S_1,\ldots,S_r$ 是整的\Term{complete}{完备}\Term{locally-noetherian}{局部 Noether} \Term{scheme}{概形}，故由 Nagata 定理（参见 EGA 0$_{\mathrm{IV}}$, 23.1.5），其\Term{normalization}{正规化} $\widetilde S_i$ 在 $S_i$ 上\Term{finite}{有限}，从而 $\widetilde S$ 在 $S$ 上\Term{finite}{有限}。于是由 Chevalley 定理（参见 EGA II, 6.7.1），若 $f\times_S\widetilde S$ 是\Term{affine}{仿射}\Term{morphism}{态射}，则 $f$ 亦然。}因此归结为 $S$ 是某个\Term{residue-field}{剩余域}\Term{algebraically-closed}{代数闭}的\Term{complete-discrete-valuation-ring}{完备离散赋值环} $A$ 的谱的情形。

\textbf{4.3.} 设 $\eta$（相应地，$s$）为 $S$ 的\Term{generic-point}{一般点}（相应地，\Term{closed-point}{闭点}）。考虑 $G_\eta$ 在 $H_\eta$ 中的像 $f_\eta(G_\eta)$。它是 $H_\eta$ 的一个闭子\Term{group-scheme}{群概形}。设 $H'$ 为 $f_\eta(G_\eta)$ 在 $H$ 中的\Term{schematic-closure}{概形闭包}。因为 $H'\to H$ \Term{affine}{仿射}（它是\Term{closed-immersion}{闭浸入}），可将 $H$ 换为 $H'$，从而假定 $H$ 在 $S$ 上\Term{flat}{平坦}且 $f_\eta$ \Term{surjective}{满射}。由于 $f_\eta$ \Term{finite}{有限}，$G,H$ 都在 $S$ 上\Term{flat}{平坦}，故有
\[\dim(G_s)=\dim(G_\eta)=\dim(H_\eta)=\dim(H_s).\]

\textbf{4.4.} 设 $H_s^0,\ldots,H_s^n$ 为 $H_s$ 的\Term{irreducible-component}{不可约分支}，其中 $H_s^0$ 表示\Term{identity-component}{单位分支}，并设 $z_0,\ldots,z_n$ 为它们的\Term{generic-point}{一般点}。因为每个\Term{local-ring}{局部环} $\mathcal O_{H,z_i}$ 的\Term{dimension}{维数} $\leq1$，\Term{morphism}{态射} $G\times_H\operatorname{Spec}(\mathcal O_{H,z_i})\to\operatorname{Spec}(\mathcal O_{H,z_i})$ 因\Term{quasi-finite}{拟有限}且\Term{separated}{分离}而\Term{affine}{仿射}（参见第 XVI 讲，引理 4.2），所以 $G$ 在 $z_i$ 的某个\Term{neighborhood}{邻域}上相对于 $H$ 是\Term{affine}{仿射}的。用 $V$ 表示 $H$ 中使 $G|_V$ 在 $V$ 上\Term{affine}{仿射}的最大\Term{open-subset}{开集}，由此可得 $V$ 包含所有 $z_i$，\footnote{编者注：我们补入了随后的句子。}
\SourcePage{277}
因此至少包含一个\Term{closed-point}{闭点} $y_i\in H_s^i$（由于 $\kappa(s)$ \Term{algebraically-closed}{代数闭}，有 $\kappa(y_i)=\kappa(s)$）。

另一方面，$V$ 显然在任意元素 $g\in G(S)$ 定义的\Term{translation}{平移}下\Term{stable}{稳定}。但 $\dim(G_s)=\dim(H_s)$，且 $f_s$ \Term{finite}{有限}，故
\[f(s):G_s^0(s)\longrightarrow H_s^0(s)\]
是\Term{surjective}{满射}。因为 $A$ \Term{complete}{完备}且 $G$ 在 $S$ 上\Term{smooth}{光滑}，\Term{canonical}{典范}映射 $G^0(S)\to G_s^0(s)$ 是\Term{surjective}{满射}；又因 $H_s^0(s)$ 在每个 $H_s^i(s)$ 上\Term{transitive-action}{传递作用}，所以 $V\supset H_s(s)$，从而（由于 $\kappa(s)$ \Term{algebraically-closed}{代数闭}）$V\supset H_s$。\footnote{编者注：当不假定 $G$ 在 $S$ 上\Term{smooth}{光滑}时，可以如下处理。取 $H_s^0$ 的一个\Term{rational-point}{有理点} $h_0$，及 $G_s^0$ 中满足 $f(g_0)=h_0$ 的一个\Term{rational-point}{有理点} $g_0$。根据第 VI$_{\mathrm B}$ 讲 5.6.1，存在\Term{commutative-diagram}{交换图}
\[
\begin{tikzcd}[ampersand replacement=\&]
S''\arrow[r,"g"]\arrow[d,"\pi"']\arrow[dr,"\phi"]\&G\arrow[d]\\
S'\arrow[r,"w"']\&S
\end{tikzcd}
\]
其中 $w$ \Term{etale}{平展}且\Term{surjective}{满射}，$\pi$ \Term{finite}{有限}且\Term{surjective}{满射}，而 $\phi^{-1}(s)$ 由唯一一个点 $s''$ 组成，且 $g(s'')=g_0$。于是\Term{morphism}{态射} $G_{S''}\to H_{S''}$ 在 $h_0y_i$ 的一个\Term{neighborhood}{邻域}上\Term{affine}{仿射}，故 $G_{S'}\to H_{S'}$ 亦然（EGA II, 6.7.1），再由\Term{faithfully-flat-descent}{忠实平坦下降}（EGA IV$_2$, 2.7.1 (xiii)），$G\to H$ 亦然。因此 $h_0y_i\in V$，由此得出 $V$ 包含 $H_s(s)$，从而包含 $H_s$。}

又因 $f_\eta$ \Term{finite}{有限}，显然 $V\supset H_\eta$，所以 $V=H$。证毕。

\SGAsection{}{参考文献}
\footnote{编者注：本讲引用的补充文献。}
\begin{description}
\item[{[BLie]}] N. Bourbaki，\emph{群与 \Term{lie-algebra}{Lie 代数}}（\emph{Groupes et algèbres de Lie}），第 I 章及第 VII--VIII 章，Hermann，1960 年及 1975 年。
\item[{[BT84]}] F. Bruhat，J. Tits，\emph{\Term{local-scheme}{局部}域上的\Term{reductive-group}{约化群} II：\Term{group-scheme}{群概形}；赋值\Term{root-datum}{根数据}的存在性}（\emph{Groupes réductifs sur un corps local II. Schémas en groupes. Existence d’une donnée radicielle valuée}），Publ. Math. I.H.É.S. \textbf{60} (1984)，5--184。
\item[{[DG70]}] M. Demazure，P. Gabriel，\emph{\Term{algebraic-group}{代数群}}（\emph{Groupes algébriques}），Masson \& North-Holland，1970。
\item[{[Ko66]}] B. Kostant，\emph{$\mathbf Z$ 上的群}（\emph{Groups over $\mathbf Z$}），90--98 页，载 \emph{\Term{algebraic-group}{代数群}及其不\Term{continuous}{连续}\Term{subgroup}{子群}}（\emph{Algebraic groups and their discontinuous subgroups}），A. Borel 与 G. D. Mostow 编，Proc. Symp. Pure Math. IX，Amer. Math. Soc.，1966。
\item[{[Lu09]}] G. Lusztig，\emph{\Term{reductive-group}{约化群}\Term{coordinate-ring}{坐标环}的一个 $\mathbf Z$-\Term{form}{形式}的研究}（\emph{Study of a $\mathbf Z$-form of the coordinate ring of a reductive group}），J. Amer. Math. Soc. \textbf{22} (2009)，第 3 期，739--769。
\item[{[Se66]}] J.-P. Serre，\emph{复\Term{semisimple-lie-algebra}{半单 Lie 代数}}（\emph{Algèbres de Lie semi-simples complexes}），Benjamin，1966。
\item[{[Sp98]}] T. A. Springer，\emph{线性\Term{algebraic-group}{代数群}}（\emph{Linear algebraic groups}），第 2 版，Birkhaüser，1998。
\item[{[Ti66]}] J. Tits，\emph{论结构常数及\Term{semisimple-lie-algebra}{半单 Lie 代数}的存在性定理}（\emph{Sur les constantes de structure et le théorème d’existence des algèbres de Lie semi-simples}），Publ. Math. I.H.É.S. \textbf{31} (1966)，21--58。
\end{description}
\SourcePage{278}
% 原 PDF 第 278 页为空白页。
